\begin{table*}[!t]
\centering
\small
\caption{Comparison with concurrent iterative-generation studies. All three examine how repeated LLM inference transforms text, but differ in chain operator, analytic lens, and downstream goal.}
\label{tab:iterative-comparison}
\resizebox{\textwidth}{!}{%
\begin{tabular}{@{}l p{2.2cm} p{2.4cm} p{2.4cm} p{2.0cm} c p{3.4cm}@{}}
\toprule
\textbf{Study} & \textbf{Chain Type} & \textbf{Feature Space} & \textbf{Analysis Goal} & \textbf{Generator Scale} & \textbf{Steps} & \textbf{Core Insight} \\
\midrule
\citet{broken_telephone}
  & Back-translation (6 bridge langs); rephrasing (ablation)
  & Surface overlap (BLEU, ROUGE, METEOR, ChrF, BERTScore) + atomic factuality (FActScore)
  & Factuality degradation rate across iterations
  & 7--9B (3 open models)
  & 100
  & Factual accuracy decays progressively; rate governed by source--bridge linguistic distance \\
\addlinespace
\citet{markovian_chains}
  & Rephrasing + round-trip translation
  & String recurrence time, distinct-output count, similarity plateau (METEOR), KL contraction
  & Convergence \& diversity dynamics of inference-time Markov chains
  & 7--8B (3 open + GPT-4o-mini)
  & 50
  & Greedy decoding collapses to 1--4 recurrent outputs; sampling sustains 10--20\,+ realizations \\
\addlinespace
\textbf{Ours}
  & Open-ended continuation (no meaning-preserving constraint)
  & Token-level NLL statistics (location, scale, shape) + lexical diversity
  & Depth classification + DDR decomposition into generator \& depth effects
  & 0.6--8B (8 models)
  & 5
  & NLL location drifts monotonically with depth while shape remains invariant, enabling per-generator depth profiling \\
\bottomrule
\end{tabular}}
\end{table*}

Table~\ref{tab:iterative-comparison} compares this work with two concurrent studies on iterative LLM generation.
All three examine how repeated inference transforms text, but differ in chain operator, analytic lens, and downstream goal.
Our continuation-chain setting permits richer depth signal and enables the location/shape decomposition that defines DDR.

\paragraph{Finding-level connections.}
\citet{broken_telephone}'s factuality degradation under back-translation chains parallels DDR's location drift: both document monotonic loss of input-specific information with iteration depth, though the former measures semantic content loss while DDR captures distributional compression in the NLL feature space.
Their finding that degradation rate scales with source--bridge linguistic distance mirrors our observation that elaboration prompts attenuate DDR (Appendix~\ref{app:prompt-type}): tighter semantic constraints slow both factuality decay and distributional drift.

\citet{markovian_chains}'s convergence dynamics under greedy decoding (collapse to 1--4 recurrent outputs within 50 steps) aligns with our greedy-decoding experiment where the DDR partition persists: shape features remain near-stationary (mean $|d|{=}0.05$) while location drift is maintained (mean $d{=}0.94$; Appendix~\ref{app:decoding-strategy}). This confirms that DDR is an empirical regularity of autoregressive generation independent of sampling stochasticity.
Their KL contraction analysis complements DDR: where they quantify \emph{how fast} the output distribution contracts, DDR decomposes \emph{what} contracts (location) from what is preserved (shape).
